\section{Identification from Order Statistics}

Let $H_0(\cdot\mid N)$ denote the distribution of the observed maximum. Under
the baseline assumptions,
\begin{equation}
H_0(b\mid N)=F_0(b)^N.
\end{equation}
For known $N$, the latent distribution is therefore identified by
\begin{equation}
F_0(v)=H_0(v\mid N)^{1/N}.
\end{equation}

Let $\widehat H_n$ be the empirical CDF of the observed maxima. The direct
inversion estimator is
\begin{equation}
\widehat F_{\mathrm{DI}}(v)=\widehat H_n(v)^{1/N}.
\end{equation}
It yields a density-free reserve estimator
\begin{equation}
\widehat r_{\mathrm{DI}}
\in\arg\max_r r\{1-\widehat F_{\mathrm{DI}}(r)\}.
\end{equation}
In computation, the objective is evaluated at the observed maxima. This direct
estimator is the first benchmark because it separates identification from the
approximation and optimization error introduced by WGAN-GP.

The transformation also makes the information problem transparent. When
$F_0(r_0)$ is below one, the probability of observing a maximum below the
reserve is $F_0(r_0)^N$, which can be small even for moderate $N$. Thus, tail
and quantile errors near the economically relevant reserve can be substantial
in finite samples despite point identification. The later semiparametric
analysis asks whether the scalar reserve can nevertheless be estimated and
inferred on without requiring uniformly accurate recovery of the entire
distribution.

These conclusions do not automatically extend to unknown competition,
correlated bidder values, auction-level unobserved heterogeneity, endogenous
entry, or data containing the transaction price rather than the highest bid.
